\lecture{16}{29. October 2025}{Boundary Layer Theory}

\begin{problem}{9.2}
  For a flow at \qty{9,0}{\m\per\s} over a smooth flat plate, estimate the maximum length of the laminar layer for air and for water as the fluid.
\end{problem}

\begin{problem}{9.8}
  Air at \qty{5}{\m\per\s}, atmospheric pressure, and \qty{20}{\celsius} flows over both sides of a flat plate that is \qty{0,8}{\m} long and \qty{0,3}{\m} wide.
  Determine the total drag force on the plate.
  Determine the total drag force when the single plate is replaced by two plates each \qty{0,4}{\m} long and \qty{0,3}{\m} wide.
  Explain why there is a difference in the total drag even though the total surface area is the same.
\end{problem}

\begin{problem}{9.7}
  For the flow of air over a flat plate, plot on the same graph the laminar boundary-layer thickness as a function of distance along the plate, up to the point of transition, for freestream speeds of \qty{1}{\m\per\s}, \qty{3}{\m\per\s} and \qty{5}{\m\per\s}.
  Draw some conclusions from your plot.
\end{problem}

